\chapter{Derivation of Interference Intensity}
\label{app:interference-derivation}

This appendix derives the interference factor in \eqref{eq:ch1-interference_intensity}.

The arbitrary polarisation lies on the great circle
$\mathbf{p}_m \in \{\mathbf{p} \in \mathcal{S}^2 : \mathbf{p} \cdot \hat{\mathbf{d}}_m = 0\}$, where $\hat{\mathbf{d}}_m$ is a unit-length direction in $\mathbb{R}^3$ expressed via the arbitrary wave-vector as $\mathbf{k}_m \coloneqq \frac{2\pi}{\lambda} \hat{\mathbf{d}}_m \in \mathcal{S}^2$. The domain of $\mathbf{p}_m$ is itself an $\mathcal{S}^1 \subset \mathcal{S}^2$ and by construction $\|\mathbf{p}_m \| = 1$.

The real electric-field superposition of $N$ monochromatic plane waves, each with angular frequency $\omega \coloneqq 2\pi \nu$ and equal amplitude $E_0$, is defined as
\begin{equation}
   \mathbf{E}_{\text{tot}}(t,\mathbf{r}) = \sum_{m=1}^N E_0 \mathbf{p}_m \cos(\mathbf{k}_m \cdot \mathbf{r} - \omega t) \ ,
\end{equation}
where $\mathbf{r}\coloneqq(x,y,z)_{\mathcal{O}} \in \mathbb{R}^3$ is the position vector for observer $\mathcal{O}$.

\noindent
The time-averaged (cycle-averaged) intensity is proportional to
\begin{equation}
   I(\mathbf{r}) \propto \langle \| \mathbf{E}_{\text{tot}} \|^2 \rangle_t  = \Bigl\langle \sum_{m,n} E_0^2 (\mathbf{p}_m \cdot \mathbf{p}_n) \cos(\mathbf{k}_m \cdot \mathbf{r} - \omega t) \cos(\mathbf{k}_n \cdot \mathbf{r} - \omega t) \Bigr\rangle_t \ .
\end{equation}

\noindent
Expanding the product of cosines using the trigonometric identity $\cos(A)\cos(B) = \frac{1}{2}[ \cos(A-B) + \cos(A+B)]$ yields
\begin{equation}
   I(\mathbf{r}) \propto \langle \| \mathbf{E}_{\text{tot}} \|^2 \rangle_t  = \Bigl\langle \sum_{m,n} E_0^2 (\mathbf{p}_m \cdot \mathbf{p}_n)
   \frac{1}{2} \bigl[ \cos((\mathbf{k}_m - \mathbf{k}_n) \cdot \mathbf{r}) + \cos((\mathbf{k}_m + \mathbf{k}_n) \cdot \mathbf{r} - 2\omega t) \bigr] \Bigr\rangle_t \ .
\end{equation}

\noindent
Under time--averaging, terms $\cos((\mathbf{k}_m + \mathbf{k}_n) \cdot \mathbf{r} - 2\omega t)$ average to zero, leaving only
\[
   \bigl\langle \cos(\mathbf{k}_m \cdot \mathbf{r} - \omega t) \cos(\mathbf{k}_n \cdot \mathbf{r} - \omega t) \bigr\rangle_t = \frac{1}{2} \cos[(\mathbf{k}_m - \mathbf{k}_n) \cdot \mathbf{r}] \ .
\]

\noindent
Separating the cross-terms:
\begin{align*}
   I(\mathbf{r})/E_0^2
   & \propto \frac{1}{2}\sum_{m,n} (\mathbf{p}_m \cdot \mathbf{p}_n)
   \cos[(\mathbf{k}_m - \mathbf{k}_n) \cdot \mathbf{r}] \\
   & = \frac{1}{2}\sum_{m=n} (\mathbf{p}_m \cdot \mathbf{p}_n)
   + \frac{1}{2}\sum_{m \neq n} (\mathbf{p}_m \cdot \mathbf{p}_n)
   \cos[(\mathbf{k}_m - \mathbf{k}_n) \cdot \mathbf{r}] \ .
\end{align*}

\noindent
Since $\mathbf{p}_m \cdot \mathbf{p}_n\vert_{n=m} = 1, \ \forall m \in [N]$ we have
\[
   \frac{1}{2}\sum_{m=n} (\mathbf{p}_m \cdot \mathbf{p}_n) = \frac{1}{2} N
\]
and for $m \neq n$, each unordered pair $\{m,n\} \in [N]^2$ appears twice in the full double sum.
Grouping them gives
\[
   \frac{1}{2}\sum_{m \neq n} (\mathbf{p}_m \cdot \mathbf{p}_n) \cos[(\mathbf{k}_m - \mathbf{k}_n) \cdot \mathbf{r}] = \sum_{m < n} (\mathbf{p}_m \cdot \mathbf{p}_n) \cos[(\mathbf{k}_m - \mathbf{k}_n) \cdot \mathbf{r}] \ .
\]

Combining the above expressions and multiplying by $2$ to match the conventional normalisation gives
\begin{equation}
   I_\text{(interf)}(\mathbf{r})/E_0^2 = N + 2\sum_{m < n} (\mathbf{p}_m \cdot \mathbf{p}_n) \cos[(\mathbf{k}_m - \mathbf{k}_n) \cdot \mathbf{r}]
\end{equation}

Defining the interference factor as $\mathcal{I}_\text{DLIP} \coloneqq I_\text{interf}/E_0^2$ we arrive at
\begin{equation}
   \mathcal{I}_{\text{DLIP}}(\mathbf{r}) = N + 2 \sum_{m < n} (\mathbf{p}_m \cdot \mathbf{p}_n) \cos[(\mathbf{k}_m - \mathbf{k}_n)\cdot \mathbf{r}] \ .
\end{equation}

This completes the derivation of the general $N$-beam interference factor.
$\square$

\subsection{Special Case: Two-Beam Interference (\(N=2\))}

In the 2--beam interference case,
the general interference factor in eq.~\eqref{eq:ch1-interference_intensity} (after dividing by $E_0^2$) collapses to a single term ($m=1$, $n=2$):
\begin{equation}
   \mathcal{I}_\text{DLIP}^{(2)}(\mathbf{r}) = 2 + 2(\mathbf{p}_1 \cdot \mathbf{p}_2) \cos[(\mathbf{k}_1 - \mathbf{k}_2) \cdot \mathbf{r}] \ .
\end{equation}

If the two beams are identically polarised ($\mathbf{p}_1 \cdot \mathbf{p}_2 = 1$), then
\begin{equation}
   \mathcal{I}_\text{DLIP}^{(2)}(\mathbf{r}) = 2[1 + \cos((\mathbf{k}_1 - \mathbf{k}_2) \cdot \mathbf{r})] \ ,
\end{equation}
which matches the two-beam result by Tsibidis et al. (eq.2.1 in \cite{Fraggelakis_2021}):
\begin{equation}
   I^{(2)}(x) \propto 1 + \cos\left(\frac{4\pi}{\lambda} x \sin\left(\frac{\theta}{2} \right) \right) \ .
\end{equation}

If the polarisations differ ($\mathbf{p}_1 \cdot \mathbf{p}_2 < 1$), the contrast of the cosine fringes is reduced by that dot-product factor.

\subsection{Special case: Four-Beam Interference (\(N=4\))}

When all $(m,n) \in [N]^2$ obey $\mathbf{p}_m \cdot \mathbf{p}_n = 1$, we recover the purely TE-mode result of
Tsibidis et al. \cite{Fraggelakis_2021} for $N=4$,
\begin{equation}
   \mathcal{I}^{(4)}_\text{DLIP}(\mathbf{r}) = 4 + 2 \sum_{m<n} \cos[(\mathbf{k}_m - \mathbf{k}_n) \cdot \mathbf{r}] \ ,
\end{equation}
after choosing \(\psi_m=\pi/2\) (pure TE polarisation) and the propagation vectors
\[
   \mathbf{k}_m
   =
   \frac{2\pi}{\lambda}
   \begin{bmatrix}
      (-1)^m \sin(\theta/2)\cos(\phi_m) \\
      (-1)^m \sin(\theta/2)\sin(\phi_m) \\
      -\cos(\theta/2)
   \end{bmatrix},
   \qquad
   \phi_m=\frac{2\pi(m-1)}{N},
   \quad m\in[N].
\]
